%% EDBT 2027 regular research-paper submission.
%% IMPORTANT: Do not change the class options, geometry, fonts, margins,
%% column spacing, footer, or other format settings supplied by EDBT.
\documentclass[sigconf,review]{acmart}

%% Copied verbatim from the official EDBT 2027 template.
\geometry{twoside=true, head=13pt,
     a4paper,
     includeheadfoot, columnsep=2pc,
     top=57pt, bottom=73pt, inner=54pt, outer=54pt,
     marginparwidth=2pc,heightrounded
     }%

\AtBeginMaketitle{%
  \input{edbt-macros}
}%

\begin{document}

%% EDBT requests title case. A regular research paper needs no title prefix.
\title{Combining Structured Pruning and Batched Model Cascades for
Semantic Query Processing}

%% EDBT 2027 uses single-anonymous review: keep author names and affiliations.
%% TODO(submission): replace every placeholder and ensure the PDF author list
%% exactly matches CMT. Every author must register in CMT and provide an ORCID.
\author{Ann Remizova}
\email{TODO@university.example}
%% \orcid{0000-0000-0000-0000}
\affiliation{%
  \institution{TODO: Institution}
  \city{TODO: City}
  \country{France}
}

\renewcommand{\shortauthors}{Remizova et al.}

\begin{abstract}
Large language models enable database systems to evaluate semantic predicates
over unstructured data, but invoking a powerful model for every tuple makes
query execution slow and resource-intensive. We present a physical execution
strategy for hybrid queries that combine structured predicates with semantic
conditions. The strategy integrates three complementary optimizations:
SQL-style structured pruning to eliminate irrelevant tuples before inference,
batch prompting to evaluate multiple candidates per invocation, and a
two-level model cascade that handles straightforward cases with a smaller
model and escalates uncertain cases to a larger model. We evaluate the approach
on ten queries over IMDb metadata and reviews, in which predicates on genre,
year, and runtime are combined with review-sentiment conditions. The benchmark
compares tuple-level semantic evaluation, batching without structured pruning,
structured pruning with row-wise cascading, and their proposed combination.
In the principal same-run comparison, structured pruning with a batched cascade
reduces mean end-to-end latency from 24.97 to 5.95 seconds and mean model calls
from 24.00 to 4.69 relative to the SUQL baseline, while achieving 0.990
precision. Its lower recall exposes the central quality--efficiency trade-off
and motivates query- and workload-aware cascade calibration. The results show
that database-aware candidate reduction and adaptive batched inference can
substantially lower the execution cost of language-model-powered queries.
\end{abstract}

\keywords{semantic query processing, large language models, query
optimization, model cascades, batch prompting, structured pruning}

\maketitle

\section{Introduction}
Hybrid analytical queries increasingly combine ordinary relational predicates
with conditions that require semantic interpretation of text. For example, a
query may restrict movies by year and runtime while asking whether each review
expresses a negative opinion. A conventional database evaluates the structured
conditions cheaply, whereas the semantic predicate may require one or more
language-model calls. Applying the latter tuple by tuple creates a costly
physical plan even when most tuples can be rejected using structured data.

This paper studies how structured pruning, batch prompting, and heterogeneous
model cascading interact in a single semantic-query plan. Our design first
extracts deterministic predicates from the user question and applies them to
the structured relation. It then groups the surviving exact-ID movie--review
pairs into prompts for a smaller model. Only uncertain decisions are coalesced
into fallback batches for a larger model.

The paper makes the following contributions:
\begin{itemize}
  \item a physical plan that composes structured candidate pruning with
        batched, cheap-to-expensive semantic inference;
  \item a controlled ten-query benchmark covering different combinations of
        structured and semantic predicates; and
  \item an empirical analysis of answer quality, latency, and cheap versus
        expensive model calls across four execution strategies.
\end{itemize}

\section{Background and Related Work}
\label{sec:related}
This section will position the method relative to semantic query languages,
LLM-powered relational operators, block-wise semantic joins, and model cascades.

%% TODO: Replace these drafting notes with a literature-backed discussion.
\paragraph{Semantic query processing.}
SUQL extends SQL with primitives that compose free-text processing with
structured data access~\cite{liu2024suql}. Our baseline uses the resulting
structured-first execution pattern but still incurs tuple-level semantic
inference after pruning.

\paragraph{Batch prompting and semantic joins.}
Batch prompting amortizes request overhead by evaluating several samples in one
model invocation~\cite{cheng2023batch}. Trummer adapts the same broad principle
to semantic joins by placing blocks from both inputs into a prompt
~\cite{trummer2025semanticjoins}. Such plans can still spend inference on
candidates that deterministic predicates would remove.

\paragraph{Model cascades.}
Confidence-based routing sends easy cases to a smaller model and ambiguous
cases to a more capable model. Semantic predicates make the routing threshold
difficult to calibrate across queries.

\section{Problem Formulation}
\label{sec:problem}
Let $M$ denote a structured movie relation and $R$ a review relation joined by
movie identifier. A hybrid query contains a structured predicate $p_s(m)$ and
a semantic predicate $p_u(r)$. The target result is
\begin{equation}
  Q = \{(m,r) \in M \Join R \mid p_s(m) \wedge p_u(r)\}.
\end{equation}
The optimizer seeks to minimize latency and model usage while maintaining
precision and recall with respect to a labeled result set.

\section{Structured Batched Cascade}
\label{sec:method}
The proposed plan has three stages.

\subsection{Structured Candidate Pruning}
Deterministic predicates over year, genre, and runtime are extracted from the
question and evaluated before semantic inference. Exact identifiers then join
the surviving movies to their reviews. This stage is conservative: an unsafe
or unsupported extraction should leave the candidate set unchanged rather than
discard possible answers.

\subsection{Batched Cheap-Model Evaluation}
Candidate pairs are grouped into schema-constrained prompts. Each output must
refer to an input identifier, preventing free-form answers from introducing
unknown tuples. Batching amortizes fixed request and prompt overhead across
several candidates.

\subsection{Uncertainty-Aware Escalation}
The smaller model assigns a decision and confidence to each candidate. Decisions
outside a calibrated confidence region are accepted or rejected immediately;
uncertain candidates are coalesced into larger fallback batches for the more
capable model. The cascade threshold controls the quality--cost trade-off.

%% TODO: Add pseudocode for the complete execution plan and define the exact
%% confidence score, threshold calibration procedure, and failure handling.

\section{Experimental Methodology}
\label{sec:methodology}
\subsection{Dataset and Workload}
The current benchmark contains ten disjoint IMDb-derived tasks. Each task has
60 unique movies, 24 candidates satisfying the structured predicates, and 12
ground-truth answers. Queries vary genre, release year or runtime, and positive
or negative review sentiment. Ground truth is the intersection of deterministic
structured filters and the source sentiment split.

\subsection{Compared Execution Strategies}
We compare (i) a SUQL baseline that prunes structured predicates and evaluates
the semantic condition with the expensive model, (ii) a batched cascade without
structured pruning, (iii) structured pruning followed by a row-wise cascade,
and (iv) the proposed structured-pruned batched cascade.

\subsection{Metrics and Repetition Policy}
We report macro precision, recall, and F1 across queries, together with wall-clock
time and the number of cheap and expensive model calls. Each question--method
configuration is repeated 11 times and numeric and quality measures are averaged.

%% TODO(submission): Document hardware, model identifiers and versions,
%% decoding settings, batch sizes, thresholds, timeout policy, software commit,
%% and statistical uncertainty. Avoid comparing rows from different runs as if
%% they were paired measurements.

\section{Results}
\label{sec:results}
Table~\ref{tab:principal-results} reports the principal same-run comparison.
The combined plan is approximately $4.2\times$ faster and uses about $80\%$
fewer model calls than SUQL. It nearly preserves precision, but the lower recall
reduces macro F1. This result supports the efficiency claim while showing that
threshold selection must be improved before claiming quality preservation.

\begin{table}[t]
  \caption{Macro results over ten IMDb queries (11 repetitions per
  question--method configuration).}
  \label{tab:principal-results}
  \centering
  \small
  \begin{tabular}{lrrrr}
    \toprule
    Method & Prec. & Recall & F1 & Time (s) \\
    \midrule
    SUQL baseline & 1.000 & 0.621 & 0.756 & 24.97 \\
    Batched cascade & 0.356 & 0.487 & 0.403 & 13.80 \\
    Pruned row cascade & 0.693 & 0.214 & 0.322 & 29.66 \\
    Pruned batched cascade & 0.990 & 0.494 & 0.607 & 5.95 \\
    \bottomrule
  \end{tabular}
\end{table}

%% TODO: Add confidence intervals, per-query results, a call-count table, and
%% ablations isolating structured pruning, batching, and cascading.

\section{Discussion and Limitations}
The benchmark is currently limited to IMDb relations and sentiment predicates,
so it does not establish generality across domains or semantic operators. The
ground truth derives from an existing sentiment split rather than independent
annotation. Local-model latency depends on the hardware and serving stack, and
call count does not by itself measure tokens, energy, or monetary cost. Finally,
the cascade sacrifices recall for some predicates; future work should calibrate
thresholds per predicate and test conservative fallback policies.

\section{Conclusion}
Structured pruning and batch prompting address different sources of waste in
semantic query execution, while a heterogeneous cascade reduces reliance on the
largest model. Combining all three yields substantial latency and call-count
reductions in the current benchmark. Broader workloads and improved calibration
are needed to retain those savings while matching the strongest baseline's
recall.

%% This exact section name and its position immediately before the references
%% are required by the EDBT 2027 call. The section is outside the page limit.
\section{Artifacts}
The implementation, benchmark-construction scripts, query specifications,
derived data, raw run outputs, and plotting scripts will be submitted as EDBT
supplemental material. The artifact README will document environment creation,
model acquisition, commands for reproducing every experiment, and the mapping
from raw outputs to each table and figure in the paper.

%% TODO(submission): Add the final supplemental ZIP name or a public repository
%% URL that does not monitor reviewer access. Verify that no secrets, local
%% absolute paths, or unnecessary model files are included.

\balance
\bibliographystyle{ACM-Reference-Format}
\bibliography{references}

\end{document}
\endinput
